\section{第一届}
\subsection*{一、计算题(本题 20 分,  每小题 5分,  共4 小题)}
\begin{enumerate}
    \item 求极限 $\displaystyle\lim_{n\rightarrow\infty}\displaystyle\sum_{k=1}^{n-1}\left(1+\frac{k}{n}\right)\sin\frac{k\pi}{n^2}$.
    \item 计算 $\displaystyle\iint\limits_{\Sigma}\frac{axdydz+(z+a)^2dxdy}{\sqrt{x^2+y^2+z^2}}$, 其中 $\Sigma$ 为下半球面 $z=-\sqrt{a^2-x^2-y^2}$ 的上侧, $a$ 为大于 0 的常数.
    \item 现要设计一个容积为 $V$ 的圆柱体容器. 已知上下两底的材料费为单位面积 $a$ 元, 而侧面的材料费为单位面积 $b$ 元. 试给出最节省的设计方案, 即高和上下底面的直径之比为何值时所需费用最少？
    \item 已知 $f(x)$ 在区间 $\left(\dfrac{1}{4},\dfrac{1}{2}\right)$ 内满足 $f'(x)=\dfrac{1}{\sin^3x+\cos^3x}$, 求 $f(x)$.
\end{enumerate}

\subsection*{二、求极限(本题10分, 第(1)小题4分, 第(2)小题6分)}
\begin{enumerate}[start=5]
    \item $\displaystyle\lim_{n\rightarrow\infty}n\left[\left(1+\dfrac{1}{n}\right)^n-e\right]$.
    \item $\displaystyle\lim_{n\rightarrow\infty}\left(\dfrac{a^{\frac{1}{n}}+b^{\frac{1}{n}}+c^{\frac{1}{n}}}{3}\right)^n$, 其中 $a>0,b>0,c>0$.
\end{enumerate}

\subsection*{三、(本题10分)}
\begin{enumerate}[start=7]
    \item 设 $f(x)$ 在点 $x=1$ 的附近有定义, 且在点 $x=1$ 处可导, $f(1)=0, f'(1)=2$, 求极限
          $\displaystyle\lim_{x\rightarrow 0}\frac{f(\sin^2 x+\cos x)}{x^2+x\tan x}.$
\end{enumerate}

\subsection*{四、(本题10分)}
\begin{enumerate}[start=8]
    \item 设函数 $f(x)$ 在 $[0,+\infty)$ 上连续, 反常积分 $\displaystyle\int_0^{+\infty}f(x)dx$ 收敛, 求
          \[
              \lim_{y\rightarrow+\infty}\frac{1}{y}\int_0^y xf(x)dx.
          \]
\end{enumerate}

\subsection*{五、(本题12分)}
\begin{enumerate}[start=9]
    \item 设函数 $f(x)$ 在 $[0,1]$ 上连续, 在 $(0,1)$ 内可微, 且 $f(0)=f(1)=0, f\left(\dfrac{1}{2}\right)=1$.证明：
          \begin{enumerate}
              \item[(1)] 存在 $\xi \in \left(\dfrac{1}{2},1\right)$ 使得 $f(\xi)=\xi$.
              \item[(2)] 存在 $\eta \in (0,\xi)$ 使得 $f'(\eta)=f(\eta)-\eta+1$.
          \end{enumerate}
\end{enumerate}

\subsection*{六、(本题14分)}
\begin{enumerate}[start=10]
    \item 设 $n > 1$ 为整数, $F(x) = \displaystyle\int_0^x e^{-t}\left(1+\frac{t}{1!}+\frac{t^2}{2!}+\cdots+\frac{t^n}{n!}\right)dt$.证明：方程 $F(x) = \dfrac{n}{2}$ 在区间 $\left(\dfrac{n}{2},n\right)$ 内至少有一个根.
\end{enumerate}

\subsection*{七、(本题12分)}
\begin{enumerate}[start=11]
    \item 是否存在 $\mathbb{R}^1$ 上的可微函数 $f(x)$ 使得 $f(f(x)) = 1 + x^2 + x^4 - x^3 - x^5$？若存在, 请给出一个例子；若不存在, 请给出证明.
\end{enumerate}

\subsection*{八、(本题12分)}
\begin{enumerate}[start=12]
    \item 设 $f(x)$ 在 $[0,+\infty)$ 上一致连续, 且对于固定的 $x \in [0,+\infty)$, 当自然数 $n \rightarrow \infty$ 时, $f(x+n) \rightarrow 0$.证明：函数序列 $\{f(x+n): n = 1,2,\cdots\}$ 在 $[0,1]$ 上一致收敛于 0.
\end{enumerate}


\newpage
\section{第二届}

\subsection*{一、计算题(本题 15 分,  每小题 5分,  共3 小题)}
\begin{enumerate}
    \item $\displaystyle\lim_{x\rightarrow 0}\left(\dfrac{\sin x}{x}\right)^{\frac{1}{1-\cos x}}$.
    \item $\displaystyle\lim_{n\rightarrow\infty}\left(\frac{1}{n+1}+\frac{1}{n+2}+\cdots+\frac{1}{n+n}\right)$.
    \item 已知 $\begin{cases}
                  x = \ln(1+e^{2t}), \\
                  y = t - \arctan e^t,
              \end{cases}$ 求 $\dfrac{\mathrm{d}^2 y}{\mathrm{d} x^2}$.
\end{enumerate}

\subsection*{二、本题 10 分}
\begin{enumerate}[start=4]
    \item 求微分方程 $(2x+y-4)\mathrm{d}x+(x+y-1)\mathrm{d}y=0$ 的通解.
\end{enumerate}

\subsection*{三、(本题 15 分)}
\begin{enumerate}[start=5]
    \item 设函数 $f(x)$ 在点 $x=0$ 的某邻域内有二阶连续导数, 且 $f(0), f'(0), f''(0)$ 均不为零.证明：存在唯一一组实数 $k_1, k_2, k_3$, 使得
          \[
              \lim_{h\rightarrow 0}\frac{k_1f(h)+k_2f(2h)+k_3f(3h)-f(0)}{h^2}=0.
          \]
\end{enumerate}

\subsection*{四、(本题 17 分)}
\begin{enumerate}[start=6]
    \item 设 $\Sigma_1: \dfrac{x^2}{a^2}+\dfrac{y^2}{b^2}+\dfrac{z^2}{c^2}=1$, 其中 $a>b>c>0$, $\Sigma_2: z^2=x^2+y^2$, $\Gamma$ 为 $\Sigma_1$ 和 $\Sigma_2$ 的交线.求椭球面 $\Sigma_1$ 在 $\Gamma$ 上各点的切平面到原点距离的最大值和最小值.
\end{enumerate}

\subsection*{五、(本题 16 分)}
\begin{enumerate}[start=7]
    \item 已知 $\Sigma$ 是空间曲线 $\left\{\begin{array}{l}
                  x^2+3y^2=1, \\
                  z=0
              \end{array}\right.$ 绕 $y$ 轴旋转而成的椭球面, $S$ 表示曲面 $\Sigma$ 的上半部分 $(z\geqslant 0)$, $\Pi$ 是椭球面 $S$ 上点 $P(x,y,z)$ 处的切平面, $\rho(x,y,z)$ 是原点到切平面 $\Pi$ 的距离, $\lambda, \mu, \nu$ 表示 $S$ 取上侧时点 $P$ 处的外法线的方向余弦.
          \begin{enumerate}
              \item[(1)] 计算 $\displaystyle\iint\limits_{S}\frac{z}{\rho(x,y,z)}\mathrm{d}S$.
              \item[(2)] 计算 $\displaystyle\iint\limits_{S}z(\lambda x+3\mu y+\nu z)\mathrm{d}S$, 其中 $S$ 取上侧.
          \end{enumerate}
\end{enumerate}

\subsection*{六、(本题 12 分)}
\begin{enumerate}[start=8]
    \item 设 $f(x)$ 是 $(-\infty,+\infty)$ 上可微的正值函数, 且满足 $|f'(x)|<mf(x)$, 其中 $0<m<1$.任取实数 $a_0$, 定义 $a_n=\ln f(a_{n-1}), n=1,2,\cdots$.证明：
          $\displaystyle\sum_{n=1}^{\infty}(a_n-a_{n-1}) \text{ 绝对收敛}.$
\end{enumerate}

\subsection*{七、(本题 15 分)}
\begin{enumerate}[start=9]
    \item 问：在区间 $[0,2]$ 上是否存在连续可微的函数 $f(x)$, 满足 $f(0)=f(2)=1, |f'(x)|\leqslant 1, \left|\displaystyle\int_0^2 f(x)\mathrm{d}x\right|\leqslant 1$？请说明理由.
\end{enumerate}


\newpage
\section{第三届}

\subsection*{一、计算题(要求写出重要步骤)(本题 30 分,  每小题 6分,  共5 小题)}
\begin{enumerate}
    \item $\displaystyle\lim_{x\rightarrow 0}\dfrac{\sin^2 x - x^2 \cos^2 x}{x^2 \sin^2 x}$.
    \item $\displaystyle\lim_{x\rightarrow +\infty}\left[\left(x^3 + \frac{x}{2} - \tan \frac{1}{x}\right)e^{\frac{1}{x}} - \sqrt{1 + x^6}\right]$.
    \item 设函数 $f(x,y)$ 有二阶连续偏导数, 满足 $f_x^2 f_{yy} - 2f_x f_y f_{xy} + f_y^2 f_{xx} = 0$, 且 $f_y \neq 0$.设 $y = y(x,z)$ 是由方程 $z = f(x,y)$ 所确定的函数, 求 $\dfrac{\partial^2 y}{\partial x^2}$.
    \item 求不定积分 $I = \displaystyle\int \left(1 + x - \frac{1}{x}\right)e^{x + \frac{1}{x}}\mathrm{d}x$.
    \item 求曲面 $x^2 + y^2 = az$ 和 $z = 2a - \sqrt{x^2 + y^2} (a > 0)$ 所围立体的表面积.
\end{enumerate}

\subsection*{二、(本题 13 分)}
\begin{enumerate}[start=6]
    \item 讨论 $\displaystyle\int_0^{+\infty} \frac{x}{\cos^2 x + x^\alpha \sin^2 x} \mathrm{d}x$ 的敛散性, 其中 $\alpha$ 是一个实常数.
\end{enumerate}

\subsection*{三、(本题 13 分)}
\begin{enumerate}[start=7]
    \item 设 $f(x)$ 在 $(-\infty, +\infty)$ 上无穷次可微, 并且满足：存在 $M > 0$, 使得 $|f^{(k)}(x)| \leqslant M (-\infty < x < +\infty), \ k = 1, 2, \cdots$, 且 $f\left(\dfrac{1}{2^n}\right) = 0 (n = 1, 2, \cdots)$.求证：在 $(-\infty, +\infty)$ 上 $f(x) \equiv 0$.
\end{enumerate}

\subsection*{四、(本题 16 分, 第(1)小题6分, 第(2)小题10分)}
\begin{enumerate}[start=8]
    \item 设 $D$ 为椭圆 $\dfrac{x^2}{a^2} + \dfrac{y^2}{b^2} \leqslant 1 \ (a > b > 0)$ 且面密度为 $\rho$ 的均质薄板, $l$ 为通过椭圆焦点 $(c, 0)$ 且垂直于薄板的直线, 其中 $c = \sqrt{a^2 - b^2}$.
          \begin{enumerate}
              \item[(1)] 求薄板 $D$ 绕直线 $l$ 旋转的转动惯量 $J$.
              \item[(2)] 对于固定的转动惯量, 讨论椭圆形薄板的面积是否有最大值和最小值.
          \end{enumerate}
\end{enumerate}

\subsection*{五、(本题 12 分)}
\begin{enumerate}[start=9]
    \item 设连续可微函数 $z = z(x, y)$ 由方程 $F(xz - y, x - yz) = 0$ 唯一确定, 其中 $F(u, v)$ 具有连续的偏导数, 试求：$I = \displaystyle\oint_L (xz^2 + 2yz) \mathrm{d}y - (2xz + yz^2) \mathrm{d}x$, 其中 $L$ 为正向单位圆周.
\end{enumerate}

\subsection*{六、(本题 16 分, 第(1)小题6分, 第(2)小题10分)}
\begin{enumerate}[start=10]
    \item
          \begin{enumerate}
              \item[(1)] 求解微分方程 $\left\{\begin{array}{l}
                            \dfrac{\mathrm{d}y}{\mathrm{d}x} - xy = x e^{x^2}, \\
                            y(0) = 1.
                        \end{array}\right.$
              \item[(2)] 设 $y = f(x)$ 为上述方程的解, 证明：
                    $\displaystyle\lim_{n\rightarrow\infty}\int_0^1 \frac{n}{n^2 x^2 + 1} f(x) \mathrm{d}x = \frac{\pi}{2}.$
          \end{enumerate}
\end{enumerate}


\newpage
\section{第四届}

\subsection*{一、计算题(本题 25 分,  每小题 5分,  共5 小题)}
\begin{enumerate}
    \item $\displaystyle\lim_{x\rightarrow 0^+}\left[\ln(x\ln a)\cdot\ln\left(\frac{\ln ax}{\ln\frac{x}{a}}\right)\right] (a>1)$.
    \item 设 $f(u,v)$ 具有连续偏导数, 且满足 $f_u(u,v)+f_v(u,v)=uv$, 求 $y(x)=e^{-2x}f(x,x)$ 所满足的一阶微分方程, 并求其通解.
    \item 求在 $[0,+\infty)$ 上的可微函数 $f(x)$, 使得 $f(x)=e^{-u(x)}$, 其中 $u(x)=\displaystyle\int_0^x f(t)\mathrm{d}t$.
    \item 计算不定积分 $\displaystyle\int x\arctan x\ln(1+x^2)\mathrm{d}x$.
    \item 过直线 $\begin{cases}
                  10x+2y-2z=27, \\
                  x+y-z=0
              \end{cases}$ 作曲面 $3x^2+y^2-z^2=27$ 的切平面, 求此切平面的方程.
\end{enumerate}

\subsection*{二、(本题 15 分)}
\begin{enumerate}[start=6]
    \item 设曲面 $\Sigma$ 为
          \[
              \Sigma: z^2=x^2+y^2, \ 1\leqslant z\leqslant 2,
          \]
          其面密度为常数 $\rho$.求在原点处的质量为 1 的质点和 $\Sigma$ 之间的引力（记引力常数为 $G$）.
\end{enumerate}

\subsection*{三、(本题 15 分)}
\begin{enumerate}[start=7]
    \item 设函数 $f(x)$ 在 $[1,+\infty)$ 上连续可导, 且
          \[
              f'(x)=\frac{1}{1+f^2(x)}\left[\sqrt{\frac{1}{x}}-\sqrt{\ln\left(1+\frac{1}{x}\right)}\right],
          \]
          证明：$\displaystyle\lim_{x\rightarrow+\infty}f(x)$ 存在.
\end{enumerate}

\subsection*{四、(本题 15 分)}
\begin{enumerate}[start=8]
    \item 设函数 $f(x)$ 在 $[-2,2]$ 上二阶可导, 且 $|f(x)|<1$, 又 $f^2(0)+[f'(0)]^2=4$.试证：在 $(-2,2)$ 内至少存在一点 $\xi$, 使得 $f(\xi)+f''(\xi)=0$.
\end{enumerate}

\subsection*{五、(本题 15 分)}
\begin{enumerate}[start=9]
    \item 求二重积分 $\displaystyle\iint\limits_{x^2+y^2\leqslant 1}|x^2+y^2-x-y|\mathrm{d}x\mathrm{d}y$.
\end{enumerate}

\subsection*{六、(本题 15 分)}
\begin{enumerate}[start=10]
    \item 设对于任何收敛于零的序列 $\{x_n\}$, 级数 $\displaystyle\sum_{n=1}^{\infty}a_nx_n$ 都是收敛的, 试证明级数 $\displaystyle\sum_{n=1}^{\infty}|a_n|$ 收敛.
\end{enumerate}


\newpage
\section{第五届}

\subsection*{一、解答题(本题 28 分,  每小题 7分,  共4 小题)}
\begin{enumerate}
    \item 计算积分 $\displaystyle\int_0^{2\pi} x \int_x^{2\pi} \frac{\sin^2 t}{t^2} \mathrm{d}t \mathrm{d}x$.
    \item 设 $f(x)$ 是 $[0,1]$ 上的连续函数, 且满足 $\displaystyle\int_0^1 f(x) \mathrm{d}x = 1$, 求一个这样的函数 $f(x)$ 使得积分 $I = \displaystyle\int_0^1 (1+x^2)f^2(x) \mathrm{d}x$ 取得最小值.
    \item 设 $F(x,y,z)$ 和 $G(x,y,z)$ 有连续偏导数, $\dfrac{\partial(F,G)}{\partial(x,z)} \neq 0$, 曲线 $\Gamma: \begin{cases} F(x,y,z) = 0, \\ G(x,y,z) = 0 \end{cases}$ 过点 $P_0(x_0,y_0,z_0)$.记 $\Gamma$ 在 $xOy$ 平面上的投影曲线为 $S$, 求 $S$ 上过点 $(x_0,y_0)$ 的切线方程.
    \item 设矩阵 $A = \begin{pmatrix}
                  1 & 2 & 1 \\
                  3 & 4 & a \\
                  1 & 2 & 2
              \end{pmatrix}$, 其中 $a$ 为常数, 矩阵 $B$ 满足关系式 $AB = A - B + E$, 其中 $E$ 为单位矩阵且 $B \neq E$.试求常数 $a$ 的值.
\end{enumerate}

\subsection*{二、(本题 12 分)}
\begin{enumerate}[start=5]
    \item 设 $f(x) \in C^4(-\infty, +\infty)$, 且 $f(x+h) = f(x) + f'(x)h + \dfrac{1}{2}f''(x+\theta h)h^2$, 其中 $\theta$ 是与 $x, h$ 无关的常数, 证明 $f$ 是不超过三次的多项式.
\end{enumerate}

\subsection*{三、(本题 12 分)}
\begin{enumerate}[start=6]
    \item 设当 $x > -1$ 时, 可微函数 $f(x)$ 满足条件 $f'(x) + f(x) - \dfrac{1}{1+x} \displaystyle\int_0^x f(t) \mathrm{d}t = 0$, 且 $f(0) = 1$, 试证：当 $x \geqslant 0$ 时, 有 $e^{-x} \leqslant f(x) \leqslant 1$ 成立.
\end{enumerate}

\subsection*{四、(本题 12 分)}
\begin{enumerate}[start=7]
    \item 设 $D = \{(x,y) | 0 \leqslant x \leqslant 1, 0 \leqslant y \leqslant 1\}$, $I = \displaystyle\iint\limits_D f(x,y) \mathrm{d}x \mathrm{d}y$, 其中函数 $f(x,y)$ 在 $D$ 上有连续二阶偏导数.若对任何 $x, y$ 有 $f(0,y) = f(x,0) = 0$, 且 $\dfrac{\partial^2 f}{\partial x \partial y} \leqslant A$, 证明 $I \leqslant \frac{A}{4}$.
\end{enumerate}

\subsection*{五、(本题 12 分)}
\begin{enumerate}[start=8]
    \item 设函数 $f(x)$ 连续可导, $P=Q=R=f((x^2+y^2)z)$, 有向曲面 $\Sigma_t$ 是圆柱体 $x^2+y^2\leqslant t^2 (t>0), 0\leqslant z\leqslant 1$ 的表面, 方向朝外.记第二型曲面积分
          \[
              I_t = \iint\limits_{\Sigma_t} P \mathrm{d}y \mathrm{d}z + Q \mathrm{d}z \mathrm{d}x + R \mathrm{d}x \mathrm{d}y,
          \]
          求极限 $\displaystyle\lim_{t\rightarrow 0^+} \dfrac{I_t}{t^4}$.
\end{enumerate}

\subsection*{六、(本题 12 分)}
\begin{enumerate}[start=9]
    \item 设 $A, B$ 为 $n$ 阶正定矩阵, 求证：$AB$ 正定的充要条件是 $AB=BA$.
\end{enumerate}

\subsection*{七、(本题 12 分)}
\begin{enumerate}[start=10]
    \item 设 $\displaystyle\sum_{n=0}^{\infty} a_n x^n$ 的收敛半径为 1, $\displaystyle\lim_{n\rightarrow\infty} n a_n = 0$, 且 $\displaystyle\lim_{x\rightarrow 1^-} \displaystyle\sum_{n=0}^{\infty} a_n x^n = A$. 证明：$\displaystyle\sum_{n=0}^{\infty} a_n$ 收敛且 $\displaystyle\sum_{n=0}^{\infty} a_n = A$.
\end{enumerate}


\newpage
\section{第六届}

\subsection*{一、填空题(本题 30 分,  每小题 5分,  共6 小题)}
\begin{enumerate}
    \item 极限 $\displaystyle\lim_{x\rightarrow\infty}\dfrac{\left(\displaystyle\int_0^x e^{u^2}\mathrm{d}u\right)^2}{\displaystyle\int_0^x e^{2u^2}\mathrm{d}u}$ 的值是 \underline{\quad \quad}.
    \item 设实数 $a \neq 0$, 微分方程
          $\begin{cases}
                  y'' - a y'^2 = 0 , \\
                  y(0) = 0, y'(0) = -1
              \end{cases} $
          的解是 \underline{\quad \quad}.
    \item 设 $A = \begin{pmatrix}
                  \lambda & 0       & 0       \\
                  0       & \lambda & 0       \\
                  -1      & 1       & \lambda
              \end{pmatrix}$, 则 $A^{50} =$ \underline{\quad \quad}.
    \item 不定积分 $I = \displaystyle\int \frac{x^2 + 1}{x^4 + 1}\mathrm{d}x$ 等于 \underline{\quad \quad}.
    \item 设曲线积分 $I = \displaystyle\oint_L \frac{x\mathrm{d}y - y\mathrm{d}x}{|x| + |y|}$, 其中 $L$ 是以 $(1,0), (0,1), (-1,0), (0,-1)$ 为顶点的正方形的边界曲线, 方向为逆时针方向, 则 $I =$ \underline{\quad \quad}.
    \item 设 $D$ 是 $xOy$ 平面上由光滑封闭曲线围成的有界闭区域, 其面积为 $A > 0$, 函数 $f(x,y)$ 在该区域上连续且 $f(x,y) > 0$.记 $J_n = \left(\dfrac{1}{A}\displaystyle\iint\limits_D f^{1/n}(x,y)\mathrm{d}\sigma\right)^n$, 则极限 $\displaystyle\lim_{n\rightarrow\infty} J_n =$ \underline{\quad \quad}.
\end{enumerate}

\subsection*{二、(本题 12 分)}
\begin{enumerate}[start=7]
    \item 设 $\vec{l}_j, j = 1, 2, \cdots, n$ 是 $xOy$ 平面上点 $P_0$ 处的 $n \geqslant 2$ 个方向向量, 相邻两个向量之间的夹角为 $\dfrac{2\pi}{n}$.又设函数 $f(x, y)$ 在点 $P_0$ 处有连续偏导数, 证明：
          $\displaystyle\sum_{j=1}^{n} \frac{\partial f(P_0)}{\partial \vec{l}_j} = 0.$
\end{enumerate}

\subsection*{三、(本题 14 分)}
\begin{enumerate}[start=8]
    \item 设 $A_1, A_2, B_1, B_2$ 均为 $n$ 阶方阵, 其中 $A_2, B_2$ 可逆.证明：存在可逆矩阵 $P, Q$ 使得 $PA_iQ = B_i (i = 1, 2)$ 成立的充要条件是 $A_1A_2^{-1}$ 和 $B_1B_2^{-1}$ 相似.
\end{enumerate}

\subsection*{四、(本题 14 分)}
\begin{enumerate}[start=9]
    \item 设 $p > 0, x_1 = \dfrac{1}{4}, x_{n+1}^p = x_n^p + x_n^{2p} (n = 1, 2, \cdots)$, 证明 $\displaystyle\sum_{n=1}^{\infty} \frac{1}{1 + x_n^p}$ 收敛并求其和.
\end{enumerate}

\subsection*{五、(本题 15 分)}
\begin{enumerate}[start=10]
    \item
          \begin{enumerate}
              \item[(1)] 将 $[-\pi, \pi]$ 上的函数 $f(x) = |x|$ 展开成 Fourier 级数, 并证明
                    $\displaystyle\sum_{k=1}^{\infty} \frac{1}{k^2} = \frac{\pi^2}{6}.$
              \item[(2)] 求积分 $I = \displaystyle\int_{0}^{+\infty} \frac{u}{1 + e^u} \mathrm{d}u$.
          \end{enumerate}
\end{enumerate}

\subsection*{六、(本题 15 分)}
\begin{enumerate}[start=11]
    \item 设 $f(x, y)$ 为 $\mathbb{R}^2$ 上的非负连续函数, 极限 $I = \displaystyle\lim_{t \to +\infty} \iint\limits_{x^2 + y^2 \leqslant t^2} f(x, y) \mathrm{d}\sigma$ 存在且有限, 则称广义积分 \\
          $\displaystyle\iint\limits_{\mathbb{R}^2} f(x, y) \mathrm{d}\sigma$ 收敛于 $I$.
          \begin{enumerate}
              \item[(1)] 设 $\displaystyle\iint\limits_{\mathbb{R}^2} f(x, y) \mathrm{d}\sigma$ 收敛于 $I$, 证明：极限 $\displaystyle\lim_{t \to +\infty} \iint\limits_{-t \leqslant x, y \leqslant t} f(x, y) \mathrm{d}\sigma$ 存在且等于 $I$.
              \item[(2)] 设 $\displaystyle\iint\limits_{\mathbb{R}^2} e^{ax^2 + 2bxy + cy^2} \mathrm{d}\sigma$ 收敛于 $I$, 其中实二次型 $ax^2 + 2bxy + cy^2$ 在正交变换下的标准形为 $\lambda_1 u^2 + \lambda_2 v^2$.证明：$\lambda_1$ 和 $\lambda_2$ 都小于 $0$.
          \end{enumerate}
\end{enumerate}

\newpage
\section{第七届}

\subsection*{一、填空题(本题 30 分,  每小题 6分,  共5 小题)}
\begin{enumerate}
    \item 微分方程 $y'' - (y')^3 = 0$ 的通解为\underline{\quad \quad}.
    \item 设 $D: 1 \leqslant x^2 + y^2 \leqslant 4$, 则积分 $I = \displaystyle\iint\limits_{D} (x + y^2) e^{-(x^2 + y^2 - 4)} \mathrm{d}x \mathrm{d}y$ 的值是\underline{\quad \quad}.
    \item 设 $f(t)$ 二阶连续可导, 且 $f(t) \neq 0$, 若
          $\begin{cases}
                  x = \displaystyle\int_{0}^{t} f(s) \mathrm{d}s, \\
                  y = f(t),
              \end{cases}$
          则 $\dfrac{\mathrm{d}^2 y}{\mathrm{d}x^2} =$ \underline{\quad \quad}.
    \item 设 $\lambda_1, \lambda_2, \cdots, \lambda_n$ 是 $n$ 阶方阵 $A$ 的特征值, $f(x)$ 是多项式, 则矩阵 $f(A)$ 的行列式的值为 \underline{\quad \quad}.
    \item 极限 $\displaystyle\lim_{n \to \infty} [n \sin(\pi n! e)]$ 的值是\underline{\quad \quad}.
\end{enumerate}

\subsection*{二、(本题 14 分)}
\begin{enumerate}[start=6]
    \item 设函数 $f(u, v)$ 在全平面上有连续的偏导数, 曲面 $S$ 由方程 $f\left(\dfrac{x-a}{z-c}, \dfrac{y-b}{z-c}\right) = 0$ 确定, 证明：该曲面上的所有切平面都经过点 $(a, b, c)$.
\end{enumerate}

\subsection*{三、(本题 14 分)}
\begin{enumerate}[start=7]
    \item 设 $f(x)$ 在 $[a, b]$ 上连续, 证明：$2 \displaystyle\int_{a}^{b} f(x) \left( \int_{x}^{b} f(t) \mathrm{d}t \right) \mathrm{d}x = \left( \int_{a}^{b} f(x) \mathrm{d}x \right)^2$.
\end{enumerate}

\subsection*{四、(本题 14 分)}
\begin{enumerate}[start=8]
    \item 设 $A$ 是 $m \times n$ 矩阵, $B$ 是 $n \times p$ 矩阵, $C$ 是 $p \times q$ 矩阵, 证明：$R(AB) + R(BC) - R(B) \leqslant R(ABC)$, 其中 $R(X)$ 表示矩阵 $X$ 的秩.
\end{enumerate}

\subsection*{五、(本题 14 分)}
\begin{enumerate}[start=9]
    \item 设 $I_n = \displaystyle\int_{0}^{\pi/4} \tan^n x \mathrm{d}x$, 其中 $n$ 为正整数.
          \begin{enumerate}
              \item[(1)] 若 $n \geqslant 2$, 计算：$I_n + I_{n-2}$.
              \item[(2)] 设 $p$ 为实数, 讨论级数 $\displaystyle\sum_{n=1}^{\infty} (-1)^n I_n^p$ 的绝对收敛性和条件收敛性.
          \end{enumerate}
\end{enumerate}

\subsection*{六、(本题 14 分)}
\begin{enumerate}[start=10]
    \item 设 $P(x, y, z)$ 和 $R(x, y, z)$ 在空间上有连续偏导数, 记上半球面 $S: z = z_0 + \sqrt{r^2 - (x - x_0)^2 - (y - y_0)^2}$, 且方向向上.若对任何点 $(x_0, y_0, z_0)$ 和 $r > 0$, 第二型曲面积分 $\displaystyle\iint\limits_{S} P \mathrm{d}y \mathrm{d}z + R \mathrm{d}x \mathrm{d}y = 0$, 证明：$\dfrac{\partial P}{\partial x} \equiv 0$.
\end{enumerate}


\newpage
\section{第八届}

\subsection*{一、填空题(本题 30 分,  每小题 6分,  共5 小题)}
\begin{enumerate}
    \item 过单叶双曲面 $\frac{x^2}{4} + \frac{y^2}{2} - 2z^2 = 1$ 与球面 $x^2 + y^2 + z^2 = 4$ 的交线且与直线
          $\begin{cases}
                  x = 0, \\
                  3y + z = 0
              \end{cases}$
          垂直的平面方程为\underline{\quad \quad}.
    \item 设可微函数 $f(x, y)$ 满足 $\dfrac{\partial f}{\partial x} = -f(x, y), f\left(0, \dfrac{\pi}{2}\right) = 1$, 且
          \[
              \lim_{n \to \infty} \left( \frac{f\left(0, y + \dfrac{1}{n}\right)}{f(0, y)} \right)^n = e^{\cot y},
          \]
          则 $f(x, y) =$\underline{\quad \quad}.
    \item 已知 $A$ 为 $n$ 阶可逆反对称矩阵, $b$ 为 $n$ 元列向量, 设 $B = \left( \begin{array}{cc} A & b \\ b^T & 0 \end{array} \right)$, 则 $\text{rank}(B) =$\underline{\quad \quad}.
    \item $\displaystyle\sum_{n=1}^{100} n^{-\frac{1}{2}}$ 的整数部分为\underline{\quad \quad}.
    \item 曲线 $L_1: y = \dfrac{1}{3}x^3 + 2x (0 \leqslant x \leqslant 1)$ 绕直线 $L_2: y = \dfrac{4}{3}x$ 旋转所生成的旋转曲面的面积为\underline{\quad \quad}.
\end{enumerate}

\subsection*{二、(本题 14 分)}
\begin{enumerate}[start=6]
    \item 设 $0 < x < \dfrac{\pi}{2}$, 证明：$\dfrac{4}{\pi^2} < \dfrac{1}{x^2} - \dfrac{1}{\tan^2 x} < \dfrac{2}{3}$.
\end{enumerate}

\subsection*{三、(本题 14 分)}
\begin{enumerate}[start=7]
    \item 设 $f(x)$ 为 $(-\infty, +\infty)$ 上连续的周期为 1 的周期函数, 且满足 $0 \leqslant f(x) \leqslant 1$ 与 $\displaystyle\int_{0}^{1} f(x) \mathrm{d}x = 1$.证明：当 $0 \leqslant x \leqslant 13$ 时, 有
          \[
              \int_{0}^{\sqrt{x}} f(t) \mathrm{d}t + \int_{0}^{\sqrt{x+27}} f(t) \mathrm{d}t + \int_{0}^{\sqrt{13-x}} f(t) \mathrm{d}t \leqslant 11,
          \]
          并给出取等号的条件.
\end{enumerate}

\subsection*{四、(本题 14 分)}
\begin{enumerate}[start=8]
    \item 设函数 $f(x, y, z)$ 在区域 $\Omega = \{(x, y, z) | x^2 + y^2 + z^2 \leqslant 1\}$ 上具有连续的二阶偏导数, 且满足 $\dfrac{\partial^2 f}{\partial x^2} + \dfrac{\partial^2 f}{\partial y^2} + \dfrac{\partial^2 f}{\partial z^2} = \sqrt{x^2 + y^2 + z^2}$.计算 $I = \displaystyle\iiint\limits_{\Omega} \left( x \frac{\partial f}{\partial x} + y \frac{\partial f}{\partial y} + z \frac{\partial f}{\partial z} \right) \mathrm{d}x \mathrm{d}y \mathrm{d}z$.
\end{enumerate}

\subsection*{五、(本题 14 分)}
\begin{enumerate}[start=9]
    \item 设 $n$ 阶方阵 $A, B$ 满足 $AB = A + B$, 证明：若存在正整数 $k$, 使 $A^k = O$（$O$ 为零矩阵）, 则行列式 $|B + 2017A| = |B|$.
\end{enumerate}

\subsection*{六、(本题 14 分)}
\begin{enumerate}[start=10]
    \item 设 $a_n = \displaystyle\sum_{k=1}^{n} \frac{1}{k} - \ln n$.
          \begin{enumerate}
              \item[(1)] 证明：极限 $\displaystyle\lim_{n \to \infty} a_n$ 存在.
              \item[(2)] 记 $\displaystyle\lim_{n \to \infty} a_n = C$, 讨论级数 $\displaystyle\sum_{n=1}^{\infty} (a_n - C)$ 的敛散性.
          \end{enumerate}
\end{enumerate}


\newpage
\section{第九届}

\subsection*{一、填空题(本题 30 分,  每小题 6分,  共5 小题)}
\begin{enumerate}
    \item 极限 $\displaystyle\lim_{x \to 0} \frac{\tan x - \sin x}{x \ln(1 + \sin^2 x)} =$ \underline{\quad \quad}.
    \item 设一平面过原点和点 $(6, -3, 2)$, 且与平面 $4x - y + 2z = 8$ 垂直, 则此平面方程为 \underline{\quad \quad}.
    \item 设函数 $f(x, y)$ 具有一阶连续偏导数, 满足 $\mathrm{d}f(x, y) = ye^y \mathrm{d}x + x(1 + y)e^y \mathrm{d}y$ 及 $f(0, 0) = 0$, 则 $f(x, y) =$ \underline{\quad \quad}.
    \item 满足 $\dfrac{\mathrm{d}u(t)}{\mathrm{d}t} = u(t) + \displaystyle\int_{0}^{1} u(t) \mathrm{d}t$ 及 $u(0) = 1$ 的可微函数 $u(t) =$ \underline{\quad \quad}.
    \item 设 $a, b, c, d$ 是互不相同的正实数, $x, y, z, w$ 是实数, 满足 $a^x = bcd, b^y = cda, c^z = dab, d^w = abc$, 则行列式
          $\begin{vmatrix}
                  -x & 1  & 1  & 1  \\
                  1  & -y & 1  & 1  \\
                  1  & 1  & -z & 1  \\
                  1  & 1  & 1  & -w
              \end{vmatrix}
              =$ \underline{\quad \quad}.
\end{enumerate}

\subsection*{二、(本题 11 分)}
\begin{enumerate}[start=6]
    \item 设函数 $f(x)$ 在区间 $(0, 1)$ 内连续, 且存在两两互异的点 $x_1, x_2, x_3, x_4 \in (0, 1)$, 使得
          \[
              \alpha = \frac{f(x_1) - f(x_2)}{x_1 - x_2} < \frac{f(x_3) - f(x_4)}{x_3 - x_4} = \beta.
          \]
          证明：对任意 $\lambda \in (\alpha, \beta)$, 存在互异的点 $x_5, x_6 \in (0, 1)$, 使得 $\lambda = \dfrac{f(x_5) - f(x_6)}{x_5 - x_6}$.
\end{enumerate}

\subsection*{三、(本题 11 分)}
\begin{enumerate}[start=7]
    \item 设函数 $f(x)$ 在区间 $[0, 1]$ 上连续且 $\displaystyle\int_{0}^{1} f(x) \mathrm{d}x \neq 0$, 证明：在区间 $[0, 1]$ 上存在三个不同的点 $x_1, x_2, x_3$, 使得
          \begin{align*}
              \frac{\pi}{8} \int_{0}^{1} f(x) \mathrm{d}x & = \left[ \frac{1}{1 + x_1^2} \int_{0}^{x_1} f(t) \mathrm{d}t + f(x_1) \arctan x_1 \right] x_3        \\
                                                          & = \left[ \frac{1}{1 + x_2^2} \int_{0}^{x_2} f(t) \mathrm{d}t + f(x_2) \arctan x_2 \right] (1 - x_3).
          \end{align*}
\end{enumerate}

\subsection*{四、(本题 12 分)}
\begin{enumerate}[start=8]
    \item 求极限：$\displaystyle\lim_{n \to \infty} \left[ \sqrt[n+1]{(n+1)!} - \sqrt[n]{n!} \right]$.
\end{enumerate}

\subsection*{五、(本题 12 分)}
\begin{enumerate}[start=9]
    \item 设 $x = (x_1, x_2, \cdots, x_n)^T \in \mathbb{R}^n$, 定义 $H(x) = \displaystyle\sum_{i=1}^{n} x_i^2 - \displaystyle\sum_{i=1}^{n-1} x_i x_{i+1}$, $n \geqslant 2$.
          \begin{enumerate}
              \item[(1)] 证明：对任一非零 $x \in \mathbb{R}^n$, $H(x) > 0$.
              \item[(2)] 求 $H(x)$ 满足条件 $x_n = 1$ 的最小值.
          \end{enumerate}
\end{enumerate}

\subsection*{六、(本题 12 分)}
\begin{enumerate}[start=10]
    \item 设函数 $f(x, y)$ 在区域 $D = \{(x, y) | x^2 + y^2 \leqslant a^2\}$ 上具有一阶连续偏导数, 且满足 $f(x, y)|_{x^2 + y^2 = a^2} = a^2$, 以及 $\max\limits_{(x, y) \in D} \left[ \left( \dfrac{\partial f}{\partial x} \right)^2 + \left( \dfrac{\partial f}{\partial y} \right)^2 \right] = a^2$, 其中 $a > 0$.证明：
          \[
              \left| \iint\limits_{D} f(x, y) \mathrm{d}x \mathrm{d}y \right| \leqslant \frac{4}{3} \pi a^4.
          \]
\end{enumerate}

\subsection*{七、(本题 12 分)}
\begin{enumerate}[start=11]
    \item 设 $0 < a_n < 1, n = 1, 2, \cdots$, 且 $\displaystyle\lim_{n \to \infty} \frac{\ln \dfrac{1}{a_n}}{\ln n} = q$（$q$ 有限或 $+\infty$）.
          \begin{enumerate}
              \item[(1)] 证明：当 $q > 1$ 时级数 $\displaystyle\sum_{n=1}^{\infty} a_n$ 收敛, 当 $q < 1$ 时级数 $\displaystyle\sum_{n=1}^{\infty} a_n$ 发散.
              \item[(2)] 讨论 $q = 1$ 时级数 $\displaystyle\sum_{n=1}^{\infty} a_n$ 的收敛性并阐述理由.
          \end{enumerate}
\end{enumerate}


\newpage
\section{第十届}

\subsection*{一、填空题(本题 30 分,  每小题 6分,  共5 小题)}
\begin{enumerate}
    \item 设函数 $y = y = \begin{cases}
                  \dfrac{\sqrt{1 - a \sin^2 x} - b}{x^2}, & x \neq 0, \\
                  2,                                      & x = 0
              \end{cases} $ 在点 $x = 0$ 处连续, 则 $a + b$ 的值为 \underline{\quad \quad}.
    \item 设 $a > 0$, 则 $\displaystyle\int_{0}^{+\infty} \frac{\ln x}{x^2 + a^2} \mathrm{d}x =$\underline{\quad \quad}.
    \item 设曲线 $L$ 是空间区域 $0 \leqslant x \leqslant 1, 0 \leqslant y \leqslant 1, 0 \leqslant z \leqslant 1$ 的表面与平面 $x + y + z = \dfrac{3}{2}$ 的交线, 则 $\displaystyle\oint_L (z^2 - y^2) \mathrm{d}x + (x^2 - z^2) \mathrm{d}y + (y^2 - x^2) \mathrm{d}z \bigg| =$\underline{\quad \quad}.
    \item 设函数 $z = z(x, y)$ 由方程 $F(x - y, z) = 0$ 确定, 其中 $F(u, v)$ 具有连续二阶偏导数, 则 $\dfrac{\partial^2 z}{\partial x \partial y} =$\underline{\quad \quad}.
    \item 已知二次型 $f(x_1, x_2, \cdots, x_n) = \displaystyle\sum_{i=1}^{n} \left( x_i - \frac{x_1 + x_2 + \cdots + x_n}{n} \right)^2$, 则 $f$ 的规范形为\underline{\quad \quad}.
\end{enumerate}

\subsection*{二、(本题 12 分)}
\begin{enumerate}[start=6]
    \item 设 $f(x)$ 在区间 $(-1, 1)$ 内三阶连续可导, 满足 $f(0) = 0, f'(0) = 1, f''(0) = 0, f'''(0) = -1$.又设数列 $\{a_n\}$ 满足 $a_1 \in (0, 1), a_{n+1} = f(a_n) (n = 1, 2, 3, \cdots)$, 严格单调减少且 $\displaystyle\lim_{n \to \infty} a_n = 0$.计算 $\displaystyle\lim_{n \to \infty} n a_n^2$.
\end{enumerate}

\subsection*{三、(本题 12 分)}
\begin{enumerate}[start=7]
    \item 设 $f(x)$ 在 $(-\infty, +\infty)$ 上具有连续导数, 且 $|f(x)| \leqslant 1, f'(x) > 0, x \in (-\infty, +\infty)$.证明：对于 $0 < \alpha < \beta$, 有 $\displaystyle\lim_{n \to \infty} \int_{\alpha}^{\beta} f'(nx - \frac{1}{x}) \mathrm{d}x = 0$.
\end{enumerate}

\subsection*{四、(本题 12 分)}
\begin{enumerate}[start=8]
    \item 计算三重积分：$\displaystyle\iiint\limits_{\Omega} \frac{\mathrm{d}x \mathrm{d}y \mathrm{d}z}{(1 + x^2 + y^2 + z^2)^2}$, 其中 $\Omega: 0 \leqslant x \leqslant 1, 0 \leqslant y \leqslant 1, 0 \leqslant z \leqslant 1$.
\end{enumerate}

\subsection*{五、(本题 12 分)}
\begin{enumerate}[start=9]
    \item 求级数 $\displaystyle\sum_{n=1}^{\infty} \frac{1}{3} \cdot \frac{2}{5} \cdot \frac{3}{7} \cdots \frac{n}{2n+1} \cdot \frac{1}{n+1}$ 的和.
\end{enumerate}

\subsection*{六、(本题 11 分)}
\begin{enumerate}[start=10]
    \item 设 $A$ 是 $n$ 阶幂零矩阵, 即满足 $A^2 = O$.证明：若 $A$ 的秩为 $r$, 且 $1 \leqslant r < \dfrac{n}{2}$, 则存在 $n$ 阶可逆矩阵 $P$, 使得 $P^{-1}AP = \left( \begin{array}{ccc} O & I_r & O \\ O & O & O \end{array} \right)$, 其中 $I_r$ 为 $r$ 阶单位矩阵.
\end{enumerate}

\subsection*{七、(本题 11 分)}
\begin{enumerate}[start=11]
    \item 设 $\{u_n\}_{n=1}^{\infty}$ 为单调递减的正实数列, $\displaystyle\lim_{n \to \infty} u_n = 0$, $\{a_n\}_{n=1}^{\infty}$ 为一实数列, 级数 $\displaystyle\sum_{n=1}^{\infty} a_n u_n$ 收敛, 证明：$\displaystyle\lim_{n \to \infty} (a_1 + a_2 + \cdots + a_n) u_n = 0$.
\end{enumerate}


\newpage
\section{第十一届}

\subsection*{一、填空题(本题 30 分,  每小题 6分,  共5 小题)}
\begin{enumerate}
    \item 极限 $\displaystyle\lim_{x \to \frac{\pi}{2}} \frac{(1 - \sqrt{\sin x})(1 - \sqrt[3]{\sin x}) \cdots (1 - \sqrt[n]{\sin x})}{(1 - \sin x)^{n-1}} =$\underline{\quad \quad}.
    \item 设函数 $y = f(x)$ 由方程 $3x - y = 2 \arctan(y - 2x)$ 所确定, 则曲线 $y = f(x)$ 在点 $P\left(1 + \dfrac{\pi}{2}, 3 + \pi\right)$ 处的切线方程为\underline{\quad \quad}.
    \item 设平面曲线 $L$ 的方程为 $Ax^2 + By^2 + Cxy + Dx + Ey + F = 0$, 且通过五个点 $P_1(-1, 0), P_2(0, -1), P_3(0, 1), P_4(2, -1)$ 和 $P_5(2, 1)$, 则 $L$ 上任意两点之间的直线距离最大值为 \underline{\quad \quad}.
    \item 设 $f(x) = (x^2 + 2x - 3)^n \arctan^2 \dfrac{x}{3}$, 其中 $n$ 为正整数, 则 $f^{(n)}(-3) =$\underline{\quad \quad}.
    \item 设函数 $f(x)$ 的导数 $f'(x)$ 在 $[0, 1]$ 上连续, $f(0) = f(1) = 0$, 且满足
          \[
              \int_{0}^{1} [f'(x)]^2 \mathrm{d}x - 8 \int_{0}^{1} f(x) \mathrm{d}x + \frac{4}{3} = 0,
          \]
          则 $f(x) =$\underline{\quad \quad}.
\end{enumerate}

\subsection*{二、(本题 12 分)}
\begin{enumerate}[start=6]
    \item 求极限：$\displaystyle\lim_{n \to \infty} \sqrt{n} \left( 1 - \sum_{k=1}^{n} \frac{1}{n + \sqrt{k}} \right)$.
\end{enumerate}

\subsection*{三、(本题 12 分)}
\begin{enumerate}[start=7]
    \item 设 $F(x_1, x_2, x_3) = \displaystyle\int_{0}^{2\pi} f(x_1 + x_3 \cos \varphi, x_2 + x_3 \sin \varphi) \mathrm{d}\varphi$, 其中 $f(u, v)$ 具有二阶连续偏导数.\\
          已知
          $\dfrac{\partial F}{\partial x_i} = \displaystyle\int_{0}^{2\pi} \frac{\partial}{\partial x_i} [f(x_1 + x_3 \cos \varphi, x_2 + x_3 \sin \varphi)] \mathrm{d}\varphi$,
          \[
              \frac{\partial^2 F}{\partial x_i^2} = \int_{0}^{2\pi} \frac{\partial^2}{\partial x_i^2} [f(x_1 + x_3 \cos \varphi, x_2 + x_3 \sin \varphi)] \mathrm{d}\varphi, \quad i = 1, 2, 3,
          \]
          试求 $x_3 \left( \dfrac{\partial^2 F}{\partial x_1^2} + \dfrac{\partial^2 F}{\partial x_2^2} - \dfrac{\partial^2 F}{\partial x_3^2} \right) - \dfrac{\partial F}{\partial x_3}$ 并要求化简.
\end{enumerate}

\subsection*{四、(本题 10 分)}
\begin{enumerate}[start=8]
    \item 设函数 $f(x)$ 在 $[0, 1]$ 上具有连续导数, 且 $\displaystyle\int_{0}^{1} f(x) \mathrm{d}x = \frac{5}{2}$, $\displaystyle\int_{0}^{1} x f(x) \mathrm{d}x = \frac{3}{2}$.证明：存在 $\xi \in (0, 1)$, 使得 $f'(\xi) = 3$.
\end{enumerate}

\subsection*{五、(本题 12 分)}
\begin{enumerate}[start=9]
    \item 设 $B_1, B_2, \cdots, B_{2021}$ 为空间 $\mathbb{R}^3$ 中半径不为零的 2021 个球, $A = (a_{ij})$ 为 2021 阶方阵, 其 $(i, j)$ 元 $a_{ij}$ 为球 $B_i$ 与 $B_j$ 相交部分的体积.证明：行列式 $|E + A| > 1$, 其中 $E$ 为单位矩阵.
\end{enumerate}

\subsection*{六、(本题 12 分)}
\begin{enumerate}[start=10]
    \item 设 $\Omega$ 是由光滑的简单封闭曲面 $\Sigma$ 围成的有界闭区域, 函数 $f(x, y, z)$ 在 $\Omega$ 上具有连续二阶偏导数, 且

          $f(x, y, z)|_{(x, y, z) \in \Sigma} = 0$.记 $\nabla f$ 为 $f(x, y, z)$ 的梯度, 并令 $\Delta f = \dfrac{\partial^2 f}{\partial x^2} + \dfrac{\partial^2 f}{\partial y^2} + \dfrac{\partial^2 f}{\partial z^2}$.证明：对任意常数 $C > 0$, 恒有
          \[
              C \iiint\limits_{\Omega} f^2 \mathrm{d}x \mathrm{d}y \mathrm{d}z + \frac{1}{C} \iiint\limits_{\Omega} (\Delta f)^2 \mathrm{d}x \mathrm{d}y \mathrm{d}z \geqslant 2 \iiint\limits_{\Omega} |\nabla f|^2 \mathrm{d}x \mathrm{d}y \mathrm{d}z.
          \]
\end{enumerate}

\subsection*{七、(本题 12 分)}
\begin{enumerate}[start=11]
    \item 设 $\{u_n\}$ 是正数列, 满足 $\dfrac{u_{n+1}}{u_n} = 1 - \dfrac{\alpha}{n} + O\left(\dfrac{1}{n^\beta}\right)$, 其中常数 $\alpha > 0, \beta > 1$.
          \begin{enumerate}
              \item[(1)] 对于 $v_n = n^\alpha u_n$, 判断级数 $\displaystyle\sum_{n=1}^{\infty} \ln \frac{v_{n+1}}{v_n}$ 的敛散性；
              \item[(2)] 讨论级数 $\displaystyle\sum_{n=1}^{\infty} u_n$ 的敛散性.
          \end{enumerate}
\end{enumerate}


\newpage
\section{第十二届}

\subsection*{一、填空题(本题 30 分,  每小题 6分,  共5 小题)}
\begin{enumerate}
    \item 极限 $\displaystyle\lim_{n \to \infty} \sum_{k=1}^{n} \frac{k}{n^2} \sin^2 \left( 1 + \frac{k}{n} \right) =$\underline{\quad \quad}.
    \item 设 $P_0(1, 1, -1)$, $P_1(2, -1, 0)$, 则函数 $u = xyz + e^{xyz}$ 在点 $P_0$ 处沿 $\overrightarrow{P_0 P_1}$ 方向的方向导数为\underline{\quad \quad}.
    \item 记空间曲线 $\Gamma:
              \begin{cases}
                  x^2 + y^2 + z^2 = a^2, \\
                  x + y + z = 0,
              \end{cases}$
          则积分 $\displaystyle\int_{\Gamma} (1 + x)^2 \mathrm{d}s =$\underline{\quad \quad}.
    \item 设矩阵 $A$ 的伴随矩阵 $A^* = \begin{pmatrix}
                  1 &    &   \\
                    & 16 &   \\
                    &    & 1
              \end{pmatrix}$, 且满足 $ABA^{-1} = BA^{-1} + 3I$,
          其中 $I$ 为单位矩阵, 则 $B =$\underline{\quad \quad}.
    \item 函数 $u = x_1 + \dfrac{x_2}{x_1} + \dfrac{x_3}{x_2} + \dfrac{2}{x_3} \ (x_i > 0, i = 1, 2, 3)$ 的所有极值点为 \underline{\quad \quad}.
\end{enumerate}

\subsection*{二、(本题 12 分)}
\begin{enumerate}[start=6]
    \item 求极限：$\displaystyle\lim_{x \to 0} \frac{\sqrt{1 + x} \cdot \sqrt[4]{1 + 2x} \cdot \sqrt[6]{1 + 3x} \cdots \sqrt[2n]{1 + nx} - 1}{3\pi \arcsin x - (x^2 + 1) \arctan^3 x}$, 其中 $n$ 为正整数.
\end{enumerate}

\subsection*{三、(本题 12 分)}
\begin{enumerate}[start=7]
    \item 求幂级数 $\displaystyle\sum_{n=1}^{\infty} \left[ 1 - n \ln \left( 1 + \frac{1}{n} \right) \right] x^n$ 的收敛域.
\end{enumerate}

\subsection*{四、(本题 12 分)}
\begin{enumerate}[start=8]
    \item 设函数 $f(x)$ 在 $[a, b]$ 上连续, 在 $(a, b)$ 内二阶可导, 且
          \[
              f(a) = f(b) = 0, \quad \int_{a}^{b} f(x) \mathrm{d}x = 0.
          \]
          \begin{enumerate}
              \item[(1)] 存在互不相同的点 $x_1, x_2 \in (a, b)$, 使得 $f'(x_i) = f(x_i), i = 1, 2$；
              \item[(2)] 存在 $\xi \in (a, b), \xi \neq x_i, i = 1, 2$, 使得 $f''(\xi) = f(\xi)$.
          \end{enumerate}
\end{enumerate}

\subsection*{五、(本题 12 分)}
\begin{enumerate}[start=9]
    \item 设 $A$ 是 $n$ 阶实对称矩阵, 证明：
          \begin{enumerate}
              \item[(1)] 存在实对称矩阵 $B$, 使得 $B^{2021} = A$, 且 $AB = BA$；
              \item[(2)] 存在一个多项式 $p(x)$, 使得 $B = p(A)$；
              \item[(3)] 矩阵 $B$ 是唯一的.
          \end{enumerate}
\end{enumerate}

\subsection*{六、(本题 12 分)}
\begin{enumerate}[start=10]
    \item 设 $A_n(x, y) = \displaystyle\sum_{k=0}^{n} x^{n-k} y^k$, 其中 $0 < x, y < 1$, 证明：
          \[
              \frac{2}{2 - x - y} \leqslant \sum_{n=0}^{\infty} \frac{A_n(x, y)}{n+1} \leqslant \frac{1}{2} \left( \frac{1}{1 - x} + \frac{1}{1 - y} \right).
          \]
\end{enumerate}

\subsection*{七、(本题 10 分)}
\begin{enumerate}[start=11]
    \item 设 $f(x), g(x)$ 是 $[0, 1] \rightarrow [0, 1]$ 的连续函数, 且 $f(x)$ 单调增加, 求证：
          \[
              \int_{0}^{1} f(g(x)) \mathrm{d}x \leqslant \int_{0}^{1} f(x) \mathrm{d}x + \int_{0}^{1} g(x) \mathrm{d}x.
          \]
\end{enumerate}


\newpage
\section{第十三届}

\subsection*{一、填空题(本题 30 分,  每小题 6分,  共5 小题)}
\begin{enumerate}
    \item 已知 $a$ 和 $b$ 均为非零向量, 且 $|b| = 1$, $a$ 和 $b$ 的夹角 $\langle a, b \rangle = \dfrac{\pi}{4}$, 则极限
          $\displaystyle\lim_{x \to 0} \frac{|a + xb| - |a|}{x} = \underline{\quad \quad}.$
    \item 极限 $\displaystyle\lim_{x \to 0} \left[ 2 - \frac{\ln(1 + x)}{x} \right]^{\frac{2}{x}} =$\underline{\quad \quad}.
    \item 积分 $\displaystyle\int_{\sqrt{2}}^{2} \frac{\mathrm{d}x}{x \sqrt{x^2 - 1}} =$\underline{\quad \quad}.
    \item 设函数 $y = y(x)$ 由参数方程 $x = \dfrac{t}{1 + t^2}, y = \dfrac{t^2}{1 + t^2}$ 确定, 则曲线 $y = y(x)$ 在点 $\left( \dfrac{\sqrt{2}}{3}, \dfrac{2}{3} \right)$ 处的曲率 $\kappa =$ \underline{\quad \quad}.
    \item 设 $D$ 是由曲线 $\sqrt{x} + \sqrt{y} = 1$ 及两坐标轴围成的平面薄片型物件, 其密度函数为 $\rho(x, y) = \sqrt{x} + 2\sqrt{y}$, 则薄片型物件 $D$ 的质量 $M =$\underline{\quad \quad}.
\end{enumerate}

\subsection*{二、(本题 12 分)}
\begin{enumerate}[start=6]
    \item 求区间 $[0, 1]$ 上的连续函数 $f(x)$, 使之满足
          \[
              f(x) = 1 + (1 - x) \int_{0}^{x} y f(y) \mathrm{d}y + x \int_{x}^{1} (1 - y) f(y) \mathrm{d}y.
          \]
\end{enumerate}

\subsection*{三、(本题 12 分)}
\begin{enumerate}[start=7]
    \item 设曲面 $\Sigma$ 是由锥面 $x = \sqrt{y^2 + z^2}$, 平面 $x = 1$, 以及球面 $x^2 + y^2 + z^2 = 4$ 围成的满足 $x \geqslant \sqrt{y^2 + z^2}$ 的空间区域的外侧表面, 计算曲面积分：
          \[
              I = \varoiint_{\Sigma} [x^2 + f(xy)] \mathrm{d}yz + [y^2 + f(xz)] \mathrm{d}z \mathrm{d}x + [z^2 + f(yz)] \mathrm{d}x \mathrm{d}y,
          \]
          其中 $f(u)$ 是具有连续导数的奇函数.
\end{enumerate}

\subsection*{四、(本题 12 分)}
\begin{enumerate}[start=8]
    \item 设 $f(x)$ 是以 $2\pi$ 为周期的周期函数, 且 $f(x) = \begin{cases}
                  x, & 0 < x < \pi,       \\
                  0, & -\pi \leq x \leq 0
              \end{cases}$ 试将函数 $f(x)$ 展开成 Fourier（傅里叶）级数, 并求级数 $\displaystyle\sum_{k=1}^{\infty} \frac{(-1)^{n-1}}{n^2}$ 之和.
\end{enumerate}

\subsection*{五、(本题 12 分)}
\begin{enumerate}[start=9]
    \item 设数列 $\{a_n\}$ 满足 $a_1 = \dfrac{\pi}{2}, a_{n+1} = a_n - \dfrac{1}{n+1} \sin a_n, n \geqslant 1$.求证：数列 $\{na_n\}$ 收敛.
\end{enumerate}

\subsection*{六、(本题 10 分)}
\begin{enumerate}[start=10]
    \item 证明：$a^b + b^a \leqslant \sqrt{a} + \sqrt{b} \leqslant a^a + b^b$, 其中 $a > 0, b > 0, a + b = 1$.
\end{enumerate}

\subsection*{七、(本题 12 分)}
\begin{enumerate}[start=11]
    \item 设 $A = (a_{ij})$ 为 $n$ 阶实矩阵, $\alpha_1, \alpha_2, \cdots, \alpha_n$ 为 $A$ 的 $n$ 个列向量, 且均不为零.证明：矩阵 $A$ 的秩满足
          \[
              r(A) \geqslant \sum_{i=1}^{n} \frac{a_{ii}^2}{\alpha_i^T \alpha_i}.
          \]
\end{enumerate}


\newpage
\section{第十四届}

\subsection*{一、填空题(本题 30 分,  每小题 6分,  共5 小题)}
\begin{enumerate}
    \item 极限 $\displaystyle\lim_{x \to 0} \frac{\arctan x - x}{x - \sin x} =$ \underline{\quad \quad}.
    \item 设 $a > 0$, 则 $\displaystyle\int_{0}^{+\infty} \frac{x^3}{e^{ax}} \mathrm{d}x =$\underline{\quad \quad}.
    \item 点 $M_0(2, 2, 2)$ 关于直线 $L: \dfrac{x - 1}{3} = \dfrac{y + 4}{2} = z - 3$ 的对称点 $M_1$ 的坐标为\underline{\quad \quad}.
    \item 二元函数 $f(x, y) = 3xy - x^3 - y^3 + 3$ 的所有极值的和等于 \underline{\quad \quad}.
    \item 幂级数 $\displaystyle\sum_{n=1}^{\infty} (-1)^n \frac{1}{n 3^n} x^n$ 的收敛域为\underline{\quad \quad}.
\end{enumerate}

\subsection*{二、(本题 10 分)}
\begin{enumerate}[start=6]
    \item 用正交变换将二次曲面的方程
          \[
              x^2 - 2y^2 - 2z^2 - 4xy + 4xz + 8yz - 27 = 0
          \]
          化为标准方程, 并说明该曲面是什么曲面.
\end{enumerate}

\subsection*{三、(本题 12 分)}
\begin{enumerate}[start=7]
    \item 设函数 $f(x), g(x)$ 在 $(-\infty, +\infty)$ 上具有二阶连续导数, $f(0) = g(0) = 1$, 且对 $xOy$ 平面上的任一简单闭曲线 $C$, 曲线积分
          \[
              \oint\limits_C [y^2 f(x) + 2ye^x - 8yg(x)] \mathrm{d}x + 2[yg(x) + f(x)] \mathrm{d}y = 0,
          \]
          求函数 $f(x), g(x)$.
\end{enumerate}

\subsection*{四、(本题 12 分)}
\begin{enumerate}[start=8]
    \item 求由 $xOz$ 平面上的曲线
          $\begin{cases}
                  (x^2 + z^2)^2 = 4(x^2 - z^2), \\
                  y = 0
              \end{cases}$
          绕 $Oz$ 轴旋转而成的曲面所包围区域的体积.
\end{enumerate}

\subsection*{五、(本题 12 分)}
\begin{enumerate}[start=9]
    \item 证明下列不等式：
          \begin{enumerate}
              \item[(1)] 设 $x \in [0, \pi], t \in [0, 1]$, 则 $\sin tx \geqslant t \sin x$；
              \item[(2)] 设 $p > 0$, 则 $\displaystyle\int_{0}^{\frac{\pi}{2}} |\sin u|^p \mathrm{d}u \geqslant \frac{\pi}{2(p + 1)}$；
              \item[(3)] 设 $x \geqslant 0, p > 0$, 则 $\displaystyle\int_{0}^{x} |\sin u|^p \mathrm{d}u \geqslant \frac{x |\sin x|^p}{p + 1}$.
          \end{enumerate}
\end{enumerate}

\subsection*{六、(本题 12 分)}
\begin{enumerate}[start=10]
    \item 设函数 $f(x)$ 在闭区间 $[a, b]$ 上具有一阶连续导数, 证明：
          \[
              \int_{a}^{b} \sqrt{1 + [f'(x)]^2} \mathrm{d}x \geqslant \sqrt{(a - b)^2 + [f(a) - f(b)]^2},
          \]
          并给出等号成立的条件.
\end{enumerate}

\subsection*{七、(本题 12 分)}
\begin{enumerate}[start=11]
    \item 证明级数 $\displaystyle\sum_{n=1}^{\infty} \ln \left( 1 + \frac{1}{2n} \right) \cdot \ln \left( 1 + \frac{1}{2n+1} \right)$ 收敛, 并求其和.
\end{enumerate}


\newpage
\section{第十五届}

\subsection*{一、填空题(本题 30 分,  每小题 6分,  共5 小题)}
\begin{enumerate}
    \item 设函数
          $f(x) = \begin{cases}
                  2e^x(x - \cos x), & x > 0,   \\
                  x^2 + 3x + a,     & x \leq 0
              \end{cases}$
          在 $(-\infty, +\infty)$ 上连续, 则 $a =$\underline{\quad \quad}.
    \item 极限 $\displaystyle\lim_{x \to 0} \frac{x(x + 2)}{\sin \pi x} =$ \underline{\quad \quad}.
    \item 设函数 $z = z(x, y)$ 是由方程 $f\left(2x - \dfrac{z}{y}, 2y - \dfrac{z}{x}\right) = 2024$ 确定的隐函数, 其中 $f(u, v)$ 具有连续偏导数, 且 $xf_u + yf_v \neq 0$, 则 $x \dfrac{\partial z}{\partial x} + y \dfrac{\partial z}{\partial y} =$\underline{\quad \quad}.
    \item 在平面 $x + y + z = 0$ 上, 与直线
          $L_1: \begin{cases}
                  x + y - 1 = 0, \\
                  x - y + z + 1 = 0
              \end{cases}$
          和
          $L_2: \begin{cases}
                  2x - y + z - 1 = 0, \\
                  x + y - z + 1 = 0
              \end{cases}$
          都相交的直线的单位方向向量为\underline{\quad \quad}.
    \item 定积分 $\displaystyle\int_{0}^{1} \frac{x^4(1 - x)^4}{1 + x^2} \mathrm{d}x$ 的值等于\underline{\quad \quad}.
\end{enumerate}

\subsection*{二、(本题 12 分)}
\begin{enumerate}
    \item 已知曲线 $L: \begin{cases}
                  x = f(t), \\
                  y = \cos t
              \end{cases}$ \quad $(0 \leq t < \dfrac{\pi}{2})$, 其中 $f(t)$ 具有连续导数, 且 $f(0) = 0$.设当 $0 < t < \frac{\pi}{2}$ 时, $f'(t) > 0$, 且曲线 $L$ 的切线与 $x$ 轴的交点到切点的距离恒等于切点与点 $(-\sin t, 0)$ 之间的距离, 求函数 $f(t)$ 的表达式.
\end{enumerate}

\subsection*{三、(本题 12 分)}
\begin{enumerate}
    \item 求极限：$\displaystyle\lim_{n \to \infty} \int_{0}^{\frac{\pi}{2}} \frac{\sin 2n\theta}{\sin \theta} \mathrm{d}\theta$.
\end{enumerate}

\subsection*{四、(本题 12 分)}
\begin{enumerate}
    \item 设 $\Sigma_1$ 是以 $(0, 4, 0)$ 为顶点且与曲面 $\Sigma_2: \dfrac{x^2}{3} + \dfrac{y^2}{4} + \dfrac{z^2}{3} = 1 \ (y > 0)$ 相切的圆锥面, 求 $\Sigma_1$ 与 $\Sigma_2$ 所围成的空间区域的体积.
\end{enumerate}

\subsection*{五、(本题 12 分)}
\begin{enumerate}
    \item 设 $n$ 阶实矩阵 $A, B$ 满足 $AB = A + B$, 且存在 $n$ 阶可逆实矩阵 $P$, 使得 $P^{-1}AP$ 为对角矩阵.证明：$P^{-1}BP$ 也为对角矩阵.
\end{enumerate}

\subsection*{六、(本题 12 分)}
\begin{enumerate}
    \item 设数列 $\{a_n\}$ 定义为 $a_0 = 0, a_1 = \frac{2}{3}$, 当 $n \geqslant 1$ 时, 满足
          \[
              (n + 1)a_{n+1} = 2a_n + (n - 1)a_{n-1},
          \]
          求幂级数 $\displaystyle\sum_{n=0}^{\infty} n a_n x^n$ 的收敛域.
\end{enumerate}

\subsection*{七、(本题 10 分)}
\begin{enumerate}
    \item
          \begin{enumerate}
              \item[(1)] 证明：对于任意的实数 $r > 0$, 存在唯一的 $t \in (\pi, 2\pi)$, 使得 $e^{-rt} - \cos t + r \sin t = 0$；
              \item[(2)] 设 (1) 中的方程所确定的隐函数为 $t = t(r)$, 证明：当 $r > 0$, 且 $\pi < t < t(r)$ 时, 恒有 $r(\sin t - t \cos t) - t \sin t > 0$.
          \end{enumerate}
\end{enumerate}


\newpage
\section{第十六届非数A}

\subsection*{一、填空题(本题 30 分,  每小题 6分,  共5 小题)}
\begin{enumerate}
    \item 设 $a \neq 0$, 函数
          $f(x) = \begin{cases}
                  \dfrac{1 - \cos \sqrt{3x}}{a x}, & x > 0,   \\[8pt]
                  b + 2 \tan^5 x,                  & x \leq 0
              \end{cases}$
          在点 $x = 0$ 处连续, 则 $ab =$\underline{\quad \quad}.
    \item 设数列 $\{a_n\}$ 满足 $a_1 = 2$, 且点 $(\sqrt{a_{n+1}}, \sqrt{a_n}) \ (n \geqslant 1)$ 都在直线 $x - y - \sqrt{2} = 0$ 上, 则 $\displaystyle\lim_{n \to \infty} \frac{a_n}{(n + 1)^2} =$\underline{\quad \quad}.
    \item 函数 $f(x, y, z) = 2 - x^2 - y^2 + 2 \arctan z$ 在点 $P(1, 2, -1)$ 处的所有方向导数中, 最大的方向导数值为\underline{\quad \quad}.
    \item 设 $L$ 是由 $|x|^3 + 3|y| = 1$ 确定的封闭曲线.第一型曲线积分 $\displaystyle\int\limits_L |x|^3 \mathrm{d}s =$\underline{\quad \quad}.
    \item $\displaystyle\lim_{n \to + \infty} \frac{(1^{2025} + 2^{2025} + \cdots + n^{2025})^{2025}}{(1^{2024} + 2^{2024} + \cdots + n^{2024})^{2026}} =$\underline{\quad \quad}.
\end{enumerate}

\subsection*{二、(本题 10 分)}
\begin{enumerate}[start=6]
    \item 设函数
          $f(x) = \begin{cases}
                  \dfrac{\ln(1 + \sin^2 x) - 6\left(\sqrt[3]{2- \cos x}  - 1\right)}{x^3}, & x \neq 0, \\[10pt]
                  0,                                                                       & x = 0
              \end{cases}$
          证明：$f(x)$ 在 $x = 0$ 处可导, 并求 $f'(0)$.
\end{enumerate}

\subsection*{三、(本题 12 分)}
\begin{enumerate}[start=7]
    \item 求微分方程 $\dfrac{yy'' - y'^2}{y^2} = \ln y$ 的通解.
\end{enumerate}

\subsection*{四、(本题 12 分)}
\begin{enumerate}[start=8]
    \item 设 $a, b, c$ 是正数, $S$ 是方向朝上的上半椭球面
          $\dfrac{x^2}{a^2} + \dfrac{y^2}{b^2} + \dfrac{z^2}{c^2} = 1 \ (z \geq 0).$
          计算 $I = \displaystyle\iint\limits_{S} xy^2 \mathrm{d}y \mathrm{d}z + yz^2 \mathrm{d}z \mathrm{d}x + zx^2 \mathrm{d}x \mathrm{d}y$.
\end{enumerate}

\subsection*{五、(本题 12 分)}
\begin{enumerate} [start=9]
    \item 设 $A, B$ 均为 $n$ 阶实对称矩阵, 证明以下结论：
          \begin{enumerate}
              \item[(1)] 若 $A$ 是正定矩阵, 则存在可逆矩阵 $P$, 使得 $P^TAP = I$ （单位矩阵）, $P^TBP = D$ （对角矩阵）；
              \item[(2)] 设 $A, B$ 均为正定矩阵, 则对任意的 $\alpha \in (0, 1)$, 有
                    \[
                        \ln(|\alpha A + (1 - \alpha)B|) \geq \alpha \ln |A| + (1 - \alpha) \ln |B|.
                    \]
          \end{enumerate}
\end{enumerate}

\subsection*{六、(本题 12 分)}
\begin{enumerate}[start=10]
    \item 设 $f(x)$ 是 $[-1, 1]$ 上的二阶可导的奇函数, 证明：存在 $\xi \in (-1, 1)$, 使得
          \[
              f''(\xi) - f'(\xi) + f(1) = 0.
          \]
\end{enumerate}

\subsection*{七、(本题 12 分)}
\begin{enumerate}[start=11]
    \item 设数列 $\{a_n\}$ 满足 $a_1 = 2t - 3$, 其中 $t > 0$ 且 $t \neq 1$,
          \[a_{n+1} = \frac{3(t^{n+1} - 1)(a_n + 1) - 2(t^n + a_n) + (a_n + 1)}{a_n + 2t^n - 1} \quad (n \in \mathbb{N}).\]
          \begin{enumerate}
              \item[(1)] 利用变换 $b_n = \dfrac{a_n + 1}{t^n - 1}$ 求数列 $\{a_n\}$ 的通项表达式；
              \item[(2)] 令 $c_n = a_n - 1 (n = 1, 2, \ldots)$, 判断级数 $\displaystyle\sum_{n=1}^{\infty} c_n$ 的收敛性.
          \end{enumerate}
\end{enumerate}


\newpage
\section{第十六届非数B}

\subsection*{一、填空题(本题 30 分,  每小题 6分,  共5 小题)}
\begin{enumerate}
    \item $\displaystyle\lim_{(x,y) \to (0,0)} \frac{\sqrt{1 + x^2 + y^2} - 1 - \frac{x^2 + y^2}{2}}{(x^2 + y^2)^2} = $\underline{\quad \quad}.
    \item 设 $a \neq 0$, 函数
          $f(x) = \begin{cases}
                  \dfrac{1 - \cos \sqrt{3x}}{a x}, & x > 0,   \\[8pt]
                  b + 2 \tan^5 x,                  & x \leq 0
              \end{cases}$
          在点 $x = 0$ 处连续, 则 $ab =$\underline{\quad \quad}.
    \item 设数列 $\{a_n\}$ 满足 $a_1 = 2$, 且点 $(\sqrt{a_{n+1}}, \sqrt{a_n}) \ (n \geq 1)$ 都在直线 $x - y - \sqrt{2} = 0$ 上, 则 $\displaystyle\lim_{n \to \infty} \frac{a_n}{(n + 1)^2} =$\underline{\quad \quad}.
    \item 设函数 $z = x e^x (1 + y e^y)$, 则 $\left. \dfrac{\partial^2 z}{\partial x \partial y} \right|_{(1,1)} =$\underline{\quad \quad}.
    \item 函数 $f(x, y, z) = x - 2y + 3z$ 满足条件 $x^2 + y^2 + z^2 = 1$ 的最大值为\underline{\quad \quad}.
\end{enumerate}

\subsection*{二、(本题 10 分)}
\begin{enumerate}[start=6]
    \item 假设函数 $f(x)$ 满足 $f(0) = 0$ 且对于 $x \geqslant 0$, 有 $f'(x) = \dfrac{1}{x^2 + f^2(x) + 1}$.
          证明：$\displaystyle\lim_{x \to +\infty} f(x)$ 存在, 且不大于 $\dfrac{\pi}{2}$.
\end{enumerate}

\subsection*{三、(本题 12 分)}
\begin{enumerate}[start=7]
    \item 设 $I = \displaystyle\int_0^{\frac{\pi}{2}} \frac{\sin^2 x \cos^2 x}{1 + \tan^3 x} dx$, $J = \displaystyle\int_0^{\frac{\pi}{2}} \frac{\sin^2 x \cos^2 x}{1 + \cot^3 x} dx$.证明：$I = J$, 并求该积分的值.
\end{enumerate}

\subsection*{四、(本题 12 分)}
\begin{enumerate}[start=8]
    \item 求微分方程 $\dfrac{yy'' - (y')^2}{y^2} = \ln y$ 的通解.
\end{enumerate}

\subsection*{五、(本题 12 分)}
\begin{enumerate}[start=9]
    \item 讨论级数
          $1 - \dfrac{1}{2^p} + \dfrac{1}{3} - \dfrac{1}{4^p} + \cdots + \dfrac{1}{2n-1} - \dfrac{1}{(2n)^p} + \cdots$
          当 $p > 1$ 时的敛散性.
\end{enumerate}

\subsection*{六、(本题 12 分)}
\begin{enumerate}[start=10]
    \item 设 $A = (a_{ij})_{n \times n}$, $B = (b_{ij})_{n \times n}$ 为 $n \times n$ 阶实正定矩阵, $\text{tr}(A) = \displaystyle\sum_{i=1}^n a_{ii}$ 为矩阵 $A$ 的迹.证明：
          \begin{enumerate}
              \item[(1)] $\text{tr}((A+B)^2) \geqslant \text{tr}(A^2) + \text{tr}(B^2)$;
              \item[(2)] $\text{tr}((A+B)^3) \geqslant \text{tr}(A^3) + \text{tr}(B^3)$.
          \end{enumerate}
\end{enumerate}

\subsection*{七、(本题 12 分)}
\begin{enumerate}[start=11]
    \item 设 $f(x)$ 是 $[-1,1]$ 上的二阶可导的奇函数, 证明：存在 $\xi \in (-1,1)$, 使得
          \[
              f''(\xi) - f'(\xi) + f(1) = 0.
          \]
\end{enumerate}
